\chapter*{Diskuze}
\addcontentsline{toc}{chapter}{Diskuze}

\section{Zhodnocení dosažených výsledků}

Aplikace LifeMeter splňuje všechny požadavky vymezené v~zadání maturitního
projektu a~v~několika oblastech zadání překračuje. Zatímco zadání specifikuje
podporu systému Android 10 a~vyšší, výsledná aplikace je funkční rovněž na
platformě iOS díky sdílené kódové základně v~jazyce TypeScript. Integrace
s~jinou aplikací, požadovaná v~zadání v~jednotném čísle, byla realizována
souběžně pro obě platformy: prostřednictvím rozhraní Health Connect na
Androidu a~Apple HealthKit na iOS. Serverová část, jejíž existence ze zadání
přímo nevyplývala, byla vyvinuta jako plnohodnotný systém zahrnující REST API
s~52 explicitně definovanými endpointy, veřejnou webovou prezentaci v~Next.js,
administrátorskou konzoli pouze pro čtení, dokumentaci generovanou ve standardu
OpenAPI a~produkční nasazení na cloudové infrastruktuře s~automatickým
průběžným nasazením.

% ------------------------------------------------------------------
\section{Zdůvodnění volby technologií}
% ------------------------------------------------------------------

\subsection{React Native a Expo}

Volba frameworku React Native namísto nativního vývoje v~jazycích
Swift/Kotlin byla motivována požadavkem na podporu obou mobilních platforem
při zachování jedné kódové základny. Alternativou byl framework Flutter,
avšak jeho ekosystém zdravotních integrací (Health Connect, HealthKit)
byl v~době zahájení vývoje méně stabilní než dostupná řešení pro React Native.
Platforma Expo doplňuje React Native o~standardizovaný build systém
a~přístup k~nativním API prostřednictvím typovaných modulů; použití samotného
React Native bez Expo by vyžadovalo ruční správu nativních projektů Xcode
a~Android Gradle, což by projekt prodražilo na čase bez přidané funkční hodnoty.

\subsection{Bun a Hono}

Pro serverovou část byl zvolen runtime Bun namísto zavedených alternativ
Node.js nebo Deno. Bun dosahuje výrazně vyšší propustnosti při zpracování
HTTP požadavků a~disponuje vestavěným správcem balíčků, testovacím
prostředím a~bundlerem, čímž redukuje závislost na nástrojích třetích stran.
Webový framework Hono byl upřednostněn před Express.js pro svůj
typově bezpečný router a~nativní integraci s~knihovnou Zod pro validaci
schémat a~generování OpenAPI dokumentace. Tato volba umožnila udržovat
dokumentaci API v~synchronizaci se zdrojovým kódem bez nutnosti
samostatných nástrojů.

\subsection{PostgreSQL a volba relační databáze}

Pro perzistenci dat byla zvolena relační databáze PostgreSQL namísto
dokumentové databáze (např. MongoDB). Datový model aplikace -- jídla,
porce, živiny, uživatelské záznamy -- je přirozeně relační
a~JOIN operace přes normalizované tabulky jsou pro tento typ dotazů
výkonnější než vyhledávání v~zanořených dokumentech. PostgreSQL navíc
poskytuje nativní podporu pro fulltextové vyhledávání (\texttt{tsvector},
GIN indexy) a~rozšíření \texttt{pg\_trgm}, které byly klíčové pro
implementaci vyhledávání v~databázi přes milion potravin.

\subsection{Better Auth a správa relací}

Autentizace uživatelů je realizována knihovnou Better Auth s~relacemi
uloženými v~\texttt{HttpOnly} cookies namísto tokenů JWT. Tato volba vychází
z~doporučení OWASP~\cite{owasp2024mobile}: token JWT uložený v~dostupném
úložišti mobilního zařízení (např. \texttt{AsyncStorage}) je zranitelný
vůči útokům XSS a~krádeži tokenu. Relace spravovaná serverem a~identifikovaná
cookie, která není přístupná JavaScriptovému kódu, tento vektor útoku
eliminuje. Nevýhodou stavových relací je náročnější horizontální škálování
(nutnost sdíleného úložiště relací), která však pro projekt tohoto rozsahu
nepředstavuje praktické omezení.

\subsection{MMKV pro offline frontu}

Pro implementaci offline fronty bylo zvoleno úložiště MMKV namísto
standardní alternativy \texttt{AsyncStorage}. MMKV je implementováno
jako sdílená paměť mapovaná na soubor; v~benchmarcích dosahuje
přibližně 10--30násobně vyšší rychlosti čtení a~zápisu oproti
\texttt{AsyncStorage}~\cite{mmkv2024}. Pro frontu požadavků, do které
jsou zapisovány i~čteny záznamy při každé síťové operaci, je tato
výkonnostní výhoda znatelná zejména u~zařízení střední třídy.

% ------------------------------------------------------------------
\section{Srovnání s existujícími řešeními}
% ------------------------------------------------------------------

Aplikace LifeMeter je srovnávána s~řešeními analyzovanými v~teoretické
části práce.

Oproti Google Fit~\cite{googlefit2024} a~Apple Health~\cite{applehealth2024},
které plní primárně roli datových agregátorů a~nenabízejí vlastní rozhraní
pro zadávání výživy nebo silových tréninků, LifeMeter poskytuje kompletní
pracovní postup od zadání dat po vizualizaci trendů v~rámci jedné aplikace.
Obě platformy jsou zároveň využity jako zdroj dat (kroky, spánek, tělesná
hmotnost) prostřednictvím jejich zdravotních rozhraní.

MyFitnessPal~\cite{myfitnesspal2024} disponuje výrazně rozsáhlejší
komunitně budovanou databází potravin a~skenováním čárových kódů v~globálním
měřítku. LifeMeter oproti tomu vychází výhradně z~ověřené vědecké databáze
USDA FoodData Central~\cite{usda2024} s~licencí CC0~1.0, čímž je zaručena
konzistentní nutriční přesnost dat bez závislosti na správnosti komunitních
příspěvků. Sledování silových tréninků je v~MyFitnessPal omezeno na odhad
kalorického výdeje; LifeMeter zaznamenává kompletní strukturu tréninku
(sady, opakování, váha, typ sady, čas odpočinku) a~sleduje individuální
výkonnostní trend pro každé cvičení.

Cronometer~\cite{cronometer2024} se od LifeMeteru liší specializací na
nutriční přesnost a~sledování mikronutrientů s~vazbou na vědecké databáze.
LifeMeter mikronutrienty sleduje rovněž, avšak jeho záběr je šířší --
integruje silový trénink a~spánkové záznamy, které Cronometer v~základní
verzi nenabízí.

Vůči aplikaci Sleep Cycle~\cite{sleepcycle2024}, která se specializuje na
analýzu spánkových fází pomocí mikrofonu nebo akcelerometru, LifeMeter
nenabízí automatickou detekci spánkových fází přes senzory zařízení;
fáze spánku jsou importovány z~platformy zdravotních dat, pokud je uživatel
zaznamenává jiným zařízením (chytré hodinky apod.). Výhodou LifeMeteru je
integrace spánkových dat s~výživovými a~tréninkovými záznamy v~jednom
přehledu.

% ------------------------------------------------------------------
\section{Omezení implementace}
% ------------------------------------------------------------------

V~průběhu vývoje byla identifikována následující omezení:

Databáze potravin USDA FoodData Central obsahuje primárně americké potraviny
a~značkové výrobky distribuvané na americkém trhu. Dostupnost evropských
nebo českých potravin a~zejména privátních značek tuzemských řetězců je
omezená. Rozšíření databáze o~lokalizovaná data by vyžadovalo integraci
dalšího datového zdroje nebo vytvoření mechanismu pro přidávání vlastních
potravin uživatelem.

Automatická detekce spánkových fází bez externího zařízení (chytrých hodinek)
implementována není. Záznamy spánkových fází jsou importovány z~platformy
zdravotních dat, jejichž přítomnost závisí na vybavení uživatele.

% ------------------------------------------------------------------
\section{Možná budoucí rozšíření}
% ------------------------------------------------------------------

Na základě zpětné vazby z~uživatelského testování a~identifikovaných
omezení jsou jako nejvýznamnější potenciální rozšíření projektu
považována:

Přidání možnosti vytvářet a~upravovat vlastní potraviny uživatelem,
čímž by bylo překonáno omezení v~pokrytí lokálních potravin databází USDA.

Implementace widgetů pro domovskou obrazovku operačního systému,
které by umožnily rychlý přehled kalorického příjmu nebo denních kroků
bez nutnosti otevírat aplikaci.

Rozšíření modulu spánku o~automatické sledování prostřednictvím senzoru
akcelerometru zařízení jako alternativu k~importu z~chytrých hodinek.

Integrace s dalšími aplikacemi pro sledování zdraví a dalších fyzických 
metrik, například fitbit, cronometer, nebo garmin, by umožnila uživatelům 
centralizovat data z různých zdrojů a získat komplexnější přehled o svém 
zdraví a kondici.

Rozšíření na platformu Wear OS nebo watchOS by umožnilo zaznamenávat
trénink přímo z~chytrých hodinek bez nutnosti interakce s~telefonem.
